\section{A sourcewise minimum for pairwise shared costs}
\label{app:pair-resource-support}
The number of queries and the structure of their fees are different
parameters. We now fix the complete four-query source of
Theorem~\ref{thm:coverage-four} and determine a small cost-structure
threshold. Throughout this appendix a resource shared by a pair is
charged once when either query is purchased. Its activation reveals no
additional answer, so both policy classes keep the same four information
actions and the same full causal-skip priority comparator.

\begin{theorem}[Sourcewise pair-resource minimum, computer-assisted]
\label{thm:pair-resource-minimum}
Fix the complete source with
$W_0=\{1,7,9\}$, $W_1=\{0,4,11,15\}$ and the physical realization in
Theorem~\ref{thm:coverage-four}. Consider fees
\begin{equation}\label{eq:pair-resource-family}
 F(S)=\sum_{i\in S}c_i+
       \sum_{e\in\binom{\{0,1,2,3\}}2}
       w_e{\bf1}\{S\cap e\ne\varnothing\},\qquad c_i,w_e\geq0.
\end{equation}
The minimum number of positive distinct pair weights permitting
$R_{\mathcal A}<R_{\mathcal P}$ on this source is exactly two.
In particular, for every pair $e$, every $c_i\geq0$ and every
$\lambda\geq0$,
\[
 F(S)=\sum_{i\in S}c_i+\lambda{\bf1}\{S\cap e\ne\varnothing\}
 \quad\Longrightarrow\quad R_{\mathcal A}=R_{\mathcal P}.
\]
Two positive weights suffice with strictly positive individual charges:
\[
 (c_0,c_1,c_2,c_3)=(4,3,2,3),\qquad
 w_{12}=w_{23}=2,
\]
all other pair weights being zero. This strictly increasing coverage fee
has exact randomized worst-input values
\[
 R_{\mathcal A}=\frac{286}{29},\qquad
 R_{\mathcal P}=\frac{79}{8},\qquad
 R_{\mathcal P}-R_{\mathcal A}=\frac{3}{232}.
\]
\end{theorem}
The minimum concerns this source and the pair groups in
\eqref{eq:pair-resource-family}. It does not give an all-source resource
lower bound or a minimum for arbitrary larger resource groups. Repeated
copies of the same pair combine into a single weight. There is no separate
global setup term in the family being minimized.

\begin{proof}
We first exclude support zero or one over its entire continuous price
domain, then give an exact support-two witness. The complete monotone-cost
policy families for this unchanged source contain $40$ adaptive and $32$
priority profiles, as in the preceding appendix. The new proof reuses that
finite completeness premise. Every priority profile occurs among the
adaptive profiles; every original policy remains correct on all seven
inputs, including inputs assigned zero mass by a lower-bound prior.

\paragraph*{Universal one-pair exclusion.}
Fix a pair $e$ and a permutation $o$ of the four individual prices. Write
\[
 c_{o(j)}=d_0+\cdots+d_j\quad(0\leq j<4),\qquad
 \lambda=d_4,\qquad d\in\mathbb R_{\geq0}^{5}.
\]
The $24$ closed price-order cones cover every nonnegative price vector,
including equal and zero prices. A purchased-set fee and every fee
difference are linear in $d$. Thus a nonnegative difference on every
generator of a cone is nonnegative throughout the cone. Generator checks
are universal linear inequalities; they are not sampled minimax games.

Use the $32$ original priority columns followed by the $40$ original
adaptive columns. Their union has the same adaptive value. The proof
certificate closes a cone by either of two arguments.

First, a full-game saddle can use an original hard input $w$ and original
priority policy $p$. If $H$ is $p$'s purchased set on $w$, verify
\[
 F(S_a(w))\geq F(H)\quad\hbox{for every original adaptive }a,
 \qquad
 F(S_p(v))\leq F(H)\quad\hbox{for every }v\in W.
\]
The two optimal values then equal $F(H)$. The certificate uses only the
four fixed templates $(w,p)=(1,P_6),(7,P_6),(9,P_{19}),(1,P_{19})$,
with original zero-based priority indices. Even at a descendant cone,
these inequalities refer to all original policies and inputs.

Second, weak dominance may remove a strategy column when another is no
dearer on every surviving row, or an objective row when another is no
cheaper on every surviving column. These operations preserve the finite
mixed minimax value. At a terminal node, each remaining column $t$ must
have an original legal priority column $p_t$ satisfying
\begin{equation}\label{eq:full-row-pair-lift}
 F(S_{p_t}(v))\leq F(S_t(v))\qquad\hbox{for all seven original inputs }v.
\end{equation}
The all-original-row condition is essential. A deleted row was dominated
on a superset of the final columns. Induction through row deletions
therefore bounds it by a terminal row on every terminal mixture. An
optimal terminal mixture extended by zero to the original game attains
the original adaptive value on all inputs. Replacing each column by
$p_t$ with the same private probability, using
\eqref{eq:full-row-pair-lift}, gives $R_{\mathcal P}\leq R_{\mathcal A}$.
The reverse inequality follows from policy inclusion. Earlier row
inequalities remain valid on the further column subsets of a descendant
cone. No claim that separate dominance runs on the two classes reach the
same fixed point is needed.

For completeness of the parameter cover, a node may split on a linear
form $f$ taking both signs on its generators. The two closed halves
$f\geq0$ and $f\leq0$ cover their parent. Given positive generator $p$
and negative generator $n$, the boundary vector
\[
 f(p)n-f(n)p
\]
has nonnegative coefficients in $p,n$ and lies on $f=0$. The nonnegative
half is generated by its retained positive and zero generators together
with these boundary vectors: cancel negative contributions against
positive contributions in any parent conic decomposition. The opposite
half follows by negating $f$. All cones are pointed, since $d\geq0$, and
have full dimension five. For the accumulated complete inequality
description, a nonzero feasible ray is extreme exactly when its active
normals have rank four. Removing only candidate rays of smaller active
rank consequently preserves every extreme ray and the full cone. The
reader reconstructs both clipped generator sets and checks full rank
five. The rank test uses exact rational Gauss--Jordan elimination,
separate from the producer's primitive-integer elimination.

The recorded proof covers all $6\cdot24=144$ roots with $164$ nodes and
$10$ splits. All $154$ leaves close: $73$ by full-game saddles and $81$
by dominance and full-row lifting. There are $5{,}533$ column deletions,
$342$ row deletions and $252$ full-original-row priority lifts.
The separate reader checks all $504$ original tree/input paths,
$223{,}769$ scalar generator inequalities and $289$ exact ranks. Every
node is reached exactly once; no root or residual leaf is omitted.
Together with the cone-cover argument, these finite certificates prove
equality for every nonnegative $c,\lambda$, including all boundary faces.
The case $\lambda=0$ also excludes pair support zero.

\paragraph*{A two-pair witness.}
For the charges in the statement, the complete fee vector in mask order is
\begin{align*}
 (F(0),\ldots,F(7))&=(0,4,5,9,6,10,9,13),\\
 (F(8),\ldots,F(15))&=(5,9,10,14,9,13,12,16).
\end{align*}
It is a positive combination of four singleton coverage indicators and
two pair indicators, hence is normalized, nonnegative, monotone and
submodular. Its least proper query marginal is two. Some incomparable
sets have submodularity slack zero; strict increase is the asserted
strictness for this sparse witness.

Reuse the five legal trees $a_1,a_2,a_3,p_3,p_4$ from the preceding
appendix and its input order $\mathcal W=(0,1,4,7,9,11,15)$. Their costs
under the new fee are
\begin{align*}
 a_1&:(9,9,9,10,9,10,10),\\
 a_2&:(9,9,9,9,10,10,10),\\
 a_3&:(16,16,12,12,12,9,9),\\
 p_3&:(13,13,12,12,13,9,9),\\
 p_4&:(16,16,13,13,12,9,9).
\end{align*}
The adaptive mixture $(a_1,a_2,a_3,p_3)$ with weights
$(13,12,3,1)/29$ has input costs
\[
 (286,286,273,286,286,286,286)/29.
\]
The priority mixture $(a_1,a_2,p_3,p_4)$ with weights $(7,7,1,1)/16$
has costs $(155,155,151,158,158,158,158)/16$. These give the two upper
bounds on all original inputs. Each mixture chooses its seed before
seeing the input; each priority seed has one legal global order.

For the matching lower bounds, use the two priors on $\mathcal W$
\[
 \pi_A=(0,1,0,3,3,11,11)/29,\qquad
 \pi_P=(0,0,0,1,1,3,3)/8.
\]
The expected-cost slacks above $286/29$ for all $40$ original adaptive
columns, multiplied by $29$, are
\[
 \scriptsize
 \begin{array}{rrrrrrrrrr}
 100 & 117 & 91 & 100 & 117 & 133 & 0 & 17 & 91 & 77 \\
 47 & 44 & 133 & 119 & 89 & 86 & 47 & 16 & 33 & 3 \\
 0 & 0 & 17 & 8 & 25 & 25 & 77 & 63 & 16 & 33 \\
 13 & 30 & 47 & 16 & 33 & 3 & 10 & 27 & 27 & 0
 \end{array}
\]
and the slacks above $79/8$ for all $32$ original priority columns,
multiplied by $8$, are
\[
 \scriptsize
 \begin{array}{rrrrrrrr}
 28 & 33 & 25 & 28 & 33 & 37 & 0 & 5 \\
 25 & 21 & 13 & 12 & 37 & 24 & 13 & 4 \\
 9 & 1 & 0 & 0 & 5 & 1 & 6 & 6 \\
 4 & 9 & 3 & 8 & 3 & 8 & 0 & 1
 \end{array}
\]
All are nonnegative. The binding adaptive indices are $6,20,21,39$;
the binding priority indices are $6,18,19,30$. These are lower bounds
against the full original classes, rather than comparisons with the
five displayed templates alone. Any private mixture has prior expectation
at least its corresponding bound, and its maximum input expectation
is at least that prior expectation. This proves both exact values and
the gap $3/232$.

\paragraph*{Separate recurrence and evidence scope.}
A second point checker reconstructs the physical promise from all
coarsenings of the two three-block priors. It reads the frozen candidate,
with a research-input read audit excluding the policy catalogue and main
result. Under $\pi_A$, it minimizes, at each impure feasible-world set,
the sum of the two successor values over all unbought queries. A pure
state with purchased set $S$ has value $\operatorname{mass}(X)F(S)$.
Under $\pi_P$, it considers all $24$ global orders, with an explicit
permanent SKIP or purchase at each position and the same remaining order
in both branches. Every impure state must continue even when its prior
mass is zero. The resulting Bayes minima are $286/29$ and $79/8$;
the adaptive recurrence has $63$ states. Their lower-bound role, together
with the all-input upper mixtures, gives another proof of the same
randomized worst-input values without reading the retained matrix.

The point, mixtures and priors were frozen before full-column and
recurrence checking. They were obtained by balancing the five displayed
profiles; no price/prior grid or LP solver was used. The first cone
attempt left $113$ leaves unresolved using dominance/lifting alone.
The retained versioned repair added the four existing saddle templates
without changing the fee family or game. Unresolved certificates were
not treated as counterexamples or as proofs. The final certificate and
separate reader close the full domain. These are internal exact checks
using standard finite-game and polyhedral tools, with actual local exit
receipts; they are not external human review or proof-kernel validation.
\end{proof}

The two-pair example has $R_{\mathcal P}/R_{\mathcal A}=2291/2288$.
Its cost reduction relative to the optimal priority value is $3/2291$,
approximately $0.13\%$. It supplies an exact structural separation,
without establishing a large practical saving. Shared-resource and
coverage cost modeling are established~\cite{harris2025}; the sourcewise
support minimum is the claim being examined here. Pandora's fixed-order
optimality theorem~\cite[Theorem~4.1 of the inspected version]{pandora}
uses independent rewards and a different utility objective. Stochastic
Boolean evaluation~\cite{hellerstein2022adaptivity} uses its own input
distribution, costs and comparator. A theorem transfer requires preserving
these contracts. The bounded primary comparison leaves exact literature
priority open.

\begin{corollary}[A sparse resource-parameter neighborhood]
\label{cor:sparse-pair-resource-robust}
For the fixed source and two-pair fee of
Theorem~\ref{thm:pair-resource-minimum}, perturb each of the four
individual charges and the two positive pair weights by at most
$1/1856$. Throughout this six-dimensional parameter box, the fee remains
strictly increasing coverage and submodular, has exactly two positive
pair groups, and satisfies
\[
 R_{\mathcal P}-R_{\mathcal A}\geq\frac{3}{464}.
\]
\end{corollary}
\begin{proof}
For fixed input and seed, changing all purchased-set fees by at most $r$
changes cost by at most $r$. Private expectation, maximum over inputs and
infimum over a fixed policy class preserve this bound. Each optimal value
is therefore $1$-Lipschitz in the fee sup norm, so the gap loses at most
$2r$. Each set activates at most six of the listed resources. A coordinate
perturbation of size $\delta=1/1856$ changes its fee by at most
$6\delta=3/928$. The resulting gap is at least
\[
 \frac{3}{232}-12\delta=\frac{3}{464}.
\]
Every resource weight remains at least $2-\delta>0$. The four positive
singleton weights give strict increase; the positive coverage
representation gives submodularity. This parameter box is restricted
to the six stated resources and asserts no strict-submodular neighborhood
in all fifteen nonempty fee coordinates.
\end{proof}
